\documentclass[aspectratio=169,10pt]{beamer}

\usetheme{Madrid}
\usecolortheme{seahorse}
\usefonttheme{professionalfonts}

\usepackage{amsmath,amssymb,amsthm}
\usepackage{mathtools}
\usepackage{listings}
\usepackage{xcolor}
\usepackage{booktabs}

%% Python style
\lstdefinestyle{python}{
  language=Python,
  basicstyle=\ttfamily\scriptsize,
  keywordstyle=\color{blue!80!black}\bfseries,
  stringstyle=\color{orange!80!black},
  commentstyle=\color{green!50!black}\itshape,
  numberstyle=\tiny\color{gray},
  numbers=left, stepnumber=1,
  frame=single, breaklines=true,
  showstringspaces=false, tabsize=4,
  backgroundcolor=\color{gray!7},
  literate={\_}{\_}1,
}

\title[MC Theory -- Ch.1]{Lectures on Monte Carlo Theory}
\subtitle{Chapter 1: Introduction}
\author{Based on Leskel\"{a} \& Vihola}
\date{}

\setbeamertemplate{navigation symbols}{}
\setbeamertemplate{footline}[frame number]

\begin{document}

% ------------------------------------------------------------------ title
\begin{frame}
  \titlepage
\end{frame}

% ------------------------------------------------------------------ TOC
\begin{frame}{Chapter Outline}
  \tableofcontents
\end{frame}

\section{1.1 About the Book}
\section{1.2 Monte Carlo: Simple Probabilistic Tasks}
\section{1.3 Quasi-Monte Carlo Methods}
\section{1.4 MC vs Quasi-MC: Comparison}
\section{1.5 Bibliographical Notes}

%% ==================================================================
%%  SECTION 1.1
%% ==================================================================
\begin{frame}[allowframebreaks]{1.1 About the Book}

  \begin{block}{What is Monte Carlo (MC)?}
    \textbf{Monte Carlo methods} are computational algorithms that use
    \emph{random (or pseudo-random) numbers} to solve mathematical and
    statistical problems: numerical integration, optimisation, simulation
    of stochastic systems, and statistical inference.
  \end{block}

  \textbf{Historical origin:} Developed in the 1940s at Los Alamos by
  von~Neumann, Ulam, and Metropolis; named after the Monte Carlo casino.

  \bigskip
  \textbf{Why randomness helps:}
  \begin{itemize}
    \item High-dimensional integrals --- grid methods cost $n^d$ (exponential in $d$).
    \item NP-hard combinatorial problems (TSP, knapsack).
    \item Complex stochastic systems without closed-form solutions.
    \item Bayesian posterior distributions.
  \end{itemize}

  \framebreak

  \textbf{Book contents:}
  \begin{enumerate}
    \item \textbf{Ch.1} Introduction: motivating examples
    \item \textbf{Ch.2} Theory of random number generators and statistical tests
    \item \textbf{Ch.3} Generating random variables
    \item \textbf{Ch.4} Output analysis: confidence intervals, estimators
    \item \textbf{Ch.5} Variance reduction (stratification, antithetic, IS, CE)
    \item \textbf{Ch.6} Markov Chain MC (Metropolis, Gibbs, CFTP)
    \item \textbf{Ch.7} Stochastic optimisation (SA, cross-entropy)
    \item \textbf{Ch.8} Simulation of queues and related processes
  \end{enumerate}

  \begin{alertblock}{Key distinction: MC vs Quasi-MC}
    \textbf{Monte Carlo} uses \emph{pseudo-random} sequences (deterministic
    but statistically random-looking).
    \textbf{Quasi-Monte Carlo (QMC)} uses carefully designed
    \emph{low-discrepancy} sequences that fill $[0,1]^d$ more evenly.
  \end{alertblock}

\end{frame}

\begin{frame}[fragile]{1.1 Pseudorandom Numbers in Python}

  \begin{block}{NumPy recommended API}
    Use \texttt{default\_rng} with an explicit seed for reproducibility.
  \end{block}

\begin{lstlisting}[style=python, caption={Initialising a PRNG (Book, Sec.~1.2 intro)}]
import numpy as np
from numpy.random import PCG64, default_rng

# Recommended: default_rng with explicit seed
rng = np.random.default_rng(seed=1234)

# Generate 10 uniform samples on [0, 1)
u = rng.random(10)

# Underlying generator can be specified
rng2 = default_rng(PCG64(seed=1234))

# Successive calls produce statistically i.i.d.-looking numbers
u1 = rng2.random(5)
u2 = rng2.random(5)   # different values, same stream
\end{lstlisting}

  \textit{Successive pseudorandom numbers mimic i.i.d.\ Uniform$[0,1]$
  random variables, but are entirely deterministic given the seed.}

\end{frame}

%% ==================================================================
%%  SECTION 1.2
%% ==================================================================
\begin{frame}[allowframebreaks]{1.2 Monte Carlo: Simple Probabilistic Tasks}

  \begin{block}{Section overview}
    Three motivating examples, each easy to \emph{simulate} but hard
    (or impossible) to solve analytically at scale:
    \begin{enumerate}
      \item Simple Symmetric Random Walk (SSRW)
      \item Travelling Salesman Problem (TSP)
      \item Computing multidimensional integrals
    \end{enumerate}
  \end{block}

  \bigskip
  \textbf{Python generator for these examples (Listing 1.1 in book):}
  \[
    \texttt{rng = np.random.default\_rng(PCG64(seed=1234))}
  \]
  The \texttt{PCG64} generator is the default in NumPy $\ge 1.17$; it passes
  all standard statistical tests and has a period of $2^{128}$.

\end{frame}

% ------------------------------------------------------------------
\begin{frame}[allowframebreaks]{1.2.1 Simple Symmetric Random Walk: Setup}

  \textbf{Motivating game (Feller, Ch.~III):}
  Players A and B bet on coin tosses. A wins $+1$ on heads, $-1$ on tails.
  Starting fortune $S_0 = 0$; consider $2N$ tosses.

  \begin{block}{Definition -- Simple Symmetric Random Walk (SSRW)}
    Let $X_1, X_2, \ldots$ be i.i.d.\ with
    $\mathbb{P}(X_j = +1) = \mathbb{P}(X_j = -1) = \tfrac{1}{2}$.
    \[
      S_0 = 0, \qquad S_n \;=\; \sum_{j=1}^{n} X_j, \quad n \ge 1.
    \]
    The sequence $S_0, S_1, \ldots$ is the \emph{simple symmetric random walk}.
  \end{block}

  \textbf{Geometric interpretation:} A line segment connects
  $(i-1, S_{i-1})$ to $(i, S_i)$. A segment lies \emph{above} the
  $x$-axis if $S_{i-1} \ge 0$ or $S_i \ge 0$.

  \framebreak

  \textbf{Four questions about the realization $(S_n)_{n=0,\ldots,2N}$:}

  \begin{description}
    \item[Q1] \textbf{Last tie:} Let
      $L_{2N} = \max\{i : S_i = 0\}$. What is its distribution?
    \item[Q2] \textbf{Winning fraction:} What is the probability
      $P(\alpha,\beta)$ (for $0 \le \alpha < \beta \le 1$) that
      player~A is winning between fractions $\alpha$ and $\beta$ of the time?
    \item[Q3] \textbf{Oscillations:} How do the fluctuations of
      $(S_n)_{0\le n\le 2N}$ grow?
    \item[Q4] \textbf{Dominance:} What is the probability that one
      player is \emph{always} winning?
  \end{description}

\end{frame}

% ------------------------------------------------------------------
\begin{frame}[allowframebreaks]{1.2.1 Key Quantities of the SSRW}

  \begin{block}{Last time at zero}
    \[
      L_{2n} \;=\; \sup\{m \le 2n : S_m = 0\}
    \]
  \end{block}

  \begin{block}{Steps spent above zero}
    \[
      L^+_{2n} \;=\; \#\{j = 1, \ldots, 2n :
        S_{j-1} \ge 0 \;\text{or}\; S_j \ge 0\}
    \]
  \end{block}

  \begin{block}{Total time player A is strictly winning}
    \[
      W_{2N} = \#\{1 \le k \le 2N : S_k > 0\}
    \]
  \end{block}

  \framebreak

  \textbf{Surprising fact:}
  Intuitively, by symmetry $L^+_{2n}/(2n) \approx 1/2$.
  The \textbf{Arcsine Law} shows the opposite: the limiting density
  is U-shaped, with most mass near $0$ and $1$. The walk tends to
  spend \emph{almost all} its time on one side.

  \begin{exampleblock}{Monte Carlo estimate of the distribution of $L^+_{2n}/(2n)$}
    Run $R$ independent replications, record
    $\bigl(L^+_{2n,(r)}/(2n)\bigr)_{r=1}^R$, and plot a histogram.
    Compare to the arcsine density $p(x) = (\pi\sqrt{x(1-x)})^{-1}$.
  \end{exampleblock}

\end{frame}

% ------------------------------------------------------------------
\begin{frame}[fragile]{1.2.1 Python: Simulating the SSRW (Listing 1.1)}

\begin{lstlisting}[style=python, caption={Simple symmetric random walk simulation}]
import numpy as np
import matplotlib.pyplot as plt
from numpy.random import default_rng

rng = default_rng(seed=31415)
N   = 1000

# +1 or -1 with equal probability
steps = 1 - 2 * rng.integers(0, high=2, size=N)
S     = np.concatenate([[0], steps.cumsum()])

plt.figure(figsize=(9, 4))
plt.plot(S, lw=0.8)
plt.axhline(0, color='k', lw=0.5, ls='--')
plt.xlabel('Step $n$')
plt.ylabel('$S_n$')
plt.title(f'Simple Symmetric Random Walk, $N={N}$')
plt.tight_layout()
plt.show()
\end{lstlisting}

  A sample realization for $N = 1000$ is shown in Book Fig.~1.1.

\end{frame}

% ------------------------------------------------------------------
\begin{frame}[allowframebreaks]{1.2.1 Law of the Iterated Logarithm}

  \begin{block}{Law of the Iterated Logarithm (LIL)}
    \[
      \mathbb{P}\!\left(
        \liminf_{n \to \infty} \frac{S_n}{\sqrt{2n\log\log n}} = -1
      \right)
      \;=\;
      \mathbb{P}\!\left(
        \limsup_{n \to \infty} \frac{S_n}{\sqrt{2n\log\log n}} = +1
      \right) = 1.
    \]
    \emph{Source: [100]; cf.\ Chapter~VIII.5 in Feller [52].}
  \end{block}

  \textbf{Three scales for $S_n$:}
  \begin{center}
  \begin{tabular}{lll}
    \toprule
    Scale & Result & Verdict \\
    \midrule
    $S_n / n$ & $\to 0$ (LLN) & too strong \\
    $S_n / \sqrt{n}$ & oscillates, amplitude $\to\infty$ & too weak \\
    $S_n / \sqrt{2n\log\log n}$ & $\limsup = +1$, $\liminf = -1$ a.s. & \textbf{exact} \\
    \bottomrule
  \end{tabular}
  \end{center}

  \framebreak

  \textbf{Book Fig.~1.2:} 500 independent trajectories of length $2^{30}$.
  The paths stay within the LIL envelope $\pm\sqrt{2n\log\log n}$ (blue
  band) and touch it infinitely often; the red curves show $\pm\sqrt{n}$.

  \vspace{0.5em}
  \textbf{Simulation insight:} Even for very long walks, the envelope is
  never permanently exceeded --- but the walk \emph{will} return to it
  eventually (almost surely).

\end{frame}

\begin{frame}[fragile]{1.2.1 Python: LIL Envelope Visualization}

\begin{lstlisting}[style=python, caption={Plot random walk trajectories with LIL envelope}]
import numpy as np
import matplotlib.pyplot as plt

rng, R, N = np.random.default_rng(0), 50, 2**18
n_arr    = np.arange(1, N + 1)
envelope = np.sqrt(2 * n_arr * np.log(np.log(n_arr + 2)))

plt.figure(figsize=(10, 4))
for _ in range(R):
    steps = 1 - 2 * rng.integers(0, 2, N)
    S = np.concatenate([[0], steps.cumsum()])
    plt.plot(S, lw=0.3, alpha=0.4, color='gray')

plt.plot(n_arr,  envelope, 'b-', lw=2, label=r'$\pm\sqrt{2n\log\log n}$')
plt.plot(n_arr, -envelope, 'b-', lw=2)
plt.plot(n_arr,  np.sqrt(n_arr), 'r--', lw=1, label=r'$\pm\sqrt{n}$')
plt.plot(n_arr, -np.sqrt(n_arr), 'r--', lw=1)
plt.xlabel('Step $n$')
plt.ylabel('$S_n$')
plt.title(f'LIL envelope -- {R} trajectories, $N=2^{{18}}$')
plt.legend()
plt.tight_layout()
plt.show()
\end{lstlisting}

\end{frame}

% ------------------------------------------------------------------
\begin{frame}[allowframebreaks]{1.2.1 The Arcsine Law}

  \begin{block}{Arcsine Law (Feller [52] p.~82; Durrett [43] p.~198)}
    For $0 \le \alpha < \beta \le 1$:
    \[
      \mathbb{P}\!\left(
        \alpha \;\le\; \frac{L^+_{2n}}{2n} \;\le\; \beta
      \right)
      \;\xrightarrow{n \to \infty}\;
      \int_{\alpha}^{\beta}
      \frac{1}{\pi\,\sqrt{x(1-x)}}\,dx.
      \tag{1.2.1}
    \]
    A similar result holds for $L_{2n}$; see Durrett [43] p.~197.
  \end{block}

  \textbf{Why ``arcsine''?}
  \[
    \int_{0}^{\beta} \frac{dx}{\pi\sqrt{x(1-x)}}
    \;=\; \frac{2}{\pi}\arcsin\!\bigl(\sqrt{\beta}\bigr).
  \]

  \framebreak

  \textbf{Shape of the arcsine density:}
  \[
    p(x) = \frac{1}{\pi\sqrt{x(1-x)}}, \quad x \in (0,1).
  \]
  \begin{itemize}
    \item \textbf{U-shaped}: density $\to\infty$ at both $x=0$ and $x=1$.
    \item \textbf{Meaning:} $L^+_{2n}/(2n)$ concentrates near $0$ or $1$,
      not near $1/2$.
    \item \textbf{Counterintuitive:} the walk spends most of its time
      \emph{entirely above} or \emph{entirely below} zero.
  \end{itemize}

  \textbf{Book Fig.~1.3:} Histograms of $L^+_{2n}/(2n)$ from
  $R = 10^4$ replications with $n = 500$, overlaid with the arcsine
  density. Empirical histogram closely matches theory.

\end{frame}

\begin{frame}[fragile]{1.2.1 Python: Verifying the Arcsine Law}

\begin{lstlisting}[style=python, caption={Arcsine law: simulation vs theory}]
import numpy as np
import matplotlib.pyplot as plt

def fraction_positive(n, R=10_000, seed=0):
    rng   = np.random.default_rng(seed)
    steps = 1 - 2 * rng.integers(0, 2, (R, 2 * n))
    S     = np.hstack([np.zeros((R, 1)), steps.cumsum(axis=1)])
    # L+: step j counted if S_{j-1} >= 0 OR S_j >= 0
    above = (S[:, :-1] >= 0) | (S[:, 1:] >= 0)
    return above.sum(axis=1) / (2 * n)

fracs = fraction_positive(n=500)

x   = np.linspace(0.01, 0.99, 400)
pdf = 1.0 / (np.pi * np.sqrt(x * (1 - x)))

plt.figure(figsize=(7, 4))
plt.hist(fracs, bins=60, density=True, alpha=0.6, label='Simulated')
plt.plot(x, pdf, 'r-', lw=2, label='Arcsine density')
plt.xlabel(r'$L^+_{2n}/(2n)$')
plt.ylabel('Density')
plt.legend()
plt.title('Arcsine Law: simulation vs theory ($n=500$, $R=10^4$)')
plt.tight_layout()
plt.show()
\end{lstlisting}

\end{frame}

% ------------------------------------------------------------------
\begin{frame}[allowframebreaks]{1.2.2 Travelling Salesman Problem}

  \begin{block}{Problem statement}
    Given $n$ cities with pairwise distances $d(c_i, c_j)$, find the
    permutation $\pi^*$ that minimises the total tour length:
    \[
      \pi^* = \arg\min_{\pi \in S_n}
      \sum_{i=1}^{n} d\!\bigl(c_{\pi(i)},\, c_{\pi(i+1)}\bigr),
      \quad c_{\pi(n+1)} \coloneqq c_{\pi(1)}.
    \]
  \end{block}

  \textbf{Computational hardness:}
  \begin{itemize}
    \item \textbf{NP-complete}: no polynomial-time algorithm is known.
    \item Search space: $n!$ tours.
      $n=10$: $3.6\times10^6$;\quad $n=20$: $2.4\times10^{18}$.
    \item Exact methods feasible only for $n \lesssim 50$.
  \end{itemize}

  \textbf{MC heuristics (explored in later chapters):}
  \begin{itemize}
    \item Random restart: sample many permutations, keep the shortest.
    \item 2-opt local search with random restarts.
    \item \textbf{Simulated Annealing} (Ch.~7): accept worse solutions
      probabilistically to escape local optima.
    \item \textbf{Cross-Entropy method} (Chs.~5, 7): fit and resample.
  \end{itemize}

  \textbf{Book Fig.~1.4:} 115 USA cities from the TSPLIB benchmark.
  MC-based heuristics find near-optimal tours in seconds.

\end{frame}

\begin{frame}[fragile]{1.2.2 Python: Random Tour Search for TSP}

\begin{lstlisting}[style=python, caption={Baseline random search for TSP}]
import numpy as np

def tour_length(cities, perm):
    shifted = np.roll(perm, -1)
    return np.linalg.norm(cities[perm] - cities[shifted], axis=1).sum()

def random_tour_search(cities, R=10_000, seed=42):
    rng = np.random.default_rng(seed)
    n   = len(cities)
    best_len, best_tour = np.inf, None
    for _ in range(R):
        perm = rng.permutation(n)
        L    = tour_length(cities, perm)
        if L < best_len:
            best_len, best_tour = L, perm.copy()
    return best_tour, best_len

rng    = np.random.default_rng(0)
cities = rng.random((20, 2))       # 20 random 2-D cities
tour, length = random_tour_search(cities, R=5_000)
print(f"Best random tour length (R=5000): {length:.4f}")
# This is a weak baseline; SA and CE methods (Ch.7) do much better.
\end{lstlisting}

\end{frame}

% ------------------------------------------------------------------
\begin{frame}[allowframebreaks]{1.2.3 Computing Integrals with Monte Carlo}

  \textbf{Goal:} Estimate $I = \mathbb{E}[f(\mathbf{U})]$ where
  $\mathbf{U} \sim \mathrm{Uniform}[0,1]^d$:
  \[
    I = \int_{[0,1]^d} f(\mathbf{u})\,d\mathbf{u}.
  \]

  \begin{block}{Monte Carlo Estimator}
    Let $\mathbf{U}_1, \ldots, \mathbf{U}_R \overset{\mathrm{i.i.d.}}{\sim}
    \mathrm{Uniform}[0,1]^d$. Define $Z_i = f(\mathbf{U}_i)$. Then:
    \[
      \hat{I}_R \;=\; \frac{1}{R}\sum_{i=1}^{R} Z_i
      \;\xrightarrow{R\to\infty}\; I
      \quad\text{(SLLN)}.
    \]
  \end{block}

  \begin{block}{Variance}
    \[
      \mathrm{Var}(\hat{I}_R) = \frac{\sigma^2}{R}, \qquad
      \sigma^2 = \mathrm{Var}(Z_1) = \mathrm{Var}(f(\mathbf{U})).
    \]
    Chebyshev bound:
    $\mathbb{P}(|\hat{I}_R - I| > \varepsilon) \le \sigma^2 / (R\varepsilon^2)$.
  \end{block}

  \framebreak

  \begin{block}{Central Limit Theorem for MC}
    \[
      \sqrt{R}\,(\hat{I}_R - I)
      \;\xrightarrow{d}\; \mathcal{N}(0,\sigma^2).
    \]
    \textbf{RMSE} $= \sigma/\sqrt{R} = \mathcal{O}(R^{-1/2})$,
    \emph{regardless of the dimension $d$}.
  \end{block}

  \textbf{Dimension-free convergence:}
  \begin{center}
  \begin{tabular}{lll}
    \toprule
    Method & Error & Cost ($\varepsilon=10^{-2}$, $d=10$) \\
    \midrule
    Grid ($n^d$ pts) & $\mathcal{O}(n^{-2/d})$ & $\sim 10^{20}$ \\
    MC               & $\mathcal{O}(R^{-1/2})$ & $\sim 10^{4}$ \\
    \bottomrule
  \end{tabular}
  \end{center}

  \textbf{For $d \ge 5$}, MC is typically preferred over grid quadrature.

  \framebreak

  \textbf{Discrepancy:} The key quantity controlling QMC error:
  \[
    D(\mathbf{x}_1, \ldots, \mathbf{x}_n)
    = \sup_{B \subseteq [0,1]^d}
      \left|\frac{\#\{i : \mathbf{x}_i \in B\}}{n} - \lambda(B)\right|,
  \]
  where the sup is over all axis-aligned boxes and $\lambda(B)$ is the volume.

  \begin{itemize}
    \item For i.i.d.\ uniform: $D \approx \mathcal{O}(n^{-1/2})$.
    \item For low-discrepancy sequences: $D = \mathcal{O}((\log n)^d / n)$.
  \end{itemize}

  \textbf{Koksma--Hlawka inequality:}
  \[
    \left|\hat{I}_n - I\right| \;\le\; V(f)\cdot D(\mathbf{x}_1,\ldots,\mathbf{x}_n),
  \]
  where $V(f)$ is the Hardy--Krause variation of $f$.

\end{frame}

\begin{frame}[fragile]{1.2.3 Python: Estimating $\pi$ with MC}

\begin{lstlisting}[style=python, caption={Classic pi estimation via MC}]
import numpy as np

def estimate_pi(R, seed=0):
    rng    = np.random.default_rng(seed)
    U      = rng.random((R, 2))         # uniform on [0,1]^2
    inside = (U**2).sum(axis=1) <= 1.0  # inside unit quarter-circle
    return 4.0 * inside.mean()          # area = pi/4 -> multiply by 4

for R in [100, 1_000, 10_000, 100_000, 1_000_000]:
    pi_hat = estimate_pi(R)
    error  = abs(pi_hat - np.pi)
    se     = np.sqrt(pi_hat * (4.0 - pi_hat) / R)
    print(f"R={R:>9,}  pi_hat={pi_hat:.5f}  "
          f"|error|={error:.5f}  SE={se:.5f}")
# Expected: error ~ sigma/sqrt(R), shrinks 10x per 100x more samples
\end{lstlisting}

  \begin{exampleblock}{Convergence rate}
    Every $10\times$ increase in $R$ reduces the error by a factor of
    $\approx\!\sqrt{10} \approx 3.16$.
  \end{exampleblock}

\end{frame}

%% ==================================================================
%%  SECTION 1.3
%% ==================================================================
\begin{frame}[allowframebreaks]{1.3 Quasi-Monte Carlo (QMC) Methods}

  \begin{block}{Idea}
    Replace i.i.d.\ random points with a \textbf{deterministic,
    well-spread} sequence that covers $[0,1]^d$ more uniformly
    than random sampling.
  \end{block}

  \begin{block}{Low-discrepancy sequences}
    Several constructive sequences achieve:
    \[
      D(\mathbf{x}_1, \ldots, \mathbf{x}_n)
      \;\le\; C_d\,\frac{(\log n)^d}{n}
      + \mathcal{O}\!\left(\frac{(\log n)^{d-1}}{n}\right).
    \]
    Available in Python via \texttt{scipy.stats.qmc}:
  \end{block}

  \begin{description}
    \item[Halton] Van der Corput sequences with different prime bases per
      dimension. Simple; mild correlations for large primes.
    \item[Sobol'] Constructed via bit operations. Excellent uniformity in
      $d \ge 2$; best for $n = 2^k$.
    \item[Faure] Permuted base-$p$ sequences. Proven theoretical bounds.
    \item[LatinHypercube] Stratified: each row/column of the grid has
      exactly one sample. Good but not strictly low-discrepancy.
  \end{description}

  \textbf{Book Fig.~1.5:} $n = 1024$ PCG64 points (left) vs Sobol' (right)
  in $[0,1]^2$. The QMC points are visibly more uniform.

  \framebreak

  \begin{block}{Golden-ratio (Steinhaus) sequence -- 1D}
    \[
      x_n = \{n\alpha\}, \qquad
      \alpha = \frac{\sqrt{5}-1}{2} \approx 0.618
      \quad\text{(golden ratio)},
    \]
    where $\{x\}$ denotes the fractional part of $x$. $\alpha$ solves
    $z^2 + z - 1 = 0$ and has the continued-fraction expansion
    $\alpha = \cfrac{1}{1+\cfrac{1}{1+\cdots}}$.
    This makes $\alpha$ the ``most irrational'' number, ensuring maximal
    spread of $\{n\alpha\}$ across $[0,1)$.
  \end{block}

  \begin{block}{Kronecker sequences}
    $x_n = \{n\beta\}$ has low discrepancy whenever the partial quotients
    in $\beta = 1/(b_1 + 1/(b_2 + \cdots))$ are bounded, $b_j \le M$.
    Multidimensional Kronecker sequences $(\{n\beta_1\},\ldots,\{n\beta_d\})$
    are believed to have low discrepancy (open problem), supported by
    numerical evidence.
  \end{block}

  \framebreak

  \textbf{QMC vs MC error rate comparison:}
  \begin{center}
  \begin{tabular}{lll}
    \toprule
    Method & Error & When it wins \\
    \midrule
    MC & $\mathcal{O}(R^{-1/2})$ & High $d$, non-smooth $f$ \\
    QMC & $\mathcal{O}((\log R)^d / R)$ & Smooth $f$, moderate $d$ \\
    \bottomrule
  \end{tabular}
  \end{center}

  \textbf{Caveat:} The constant $C_d$ grows with $d$, so for very large
  $d$ the QMC advantage may not appear unless $n$ is huge.

\end{frame}

\begin{frame}[fragile]{1.3 Python: QMC Sequences with \texttt{scipy.stats.qmc}}

\begin{lstlisting}[style=python, caption={Comparing MC and QMC for pi estimation}]
import numpy as np
import matplotlib.pyplot as plt
from scipy.stats.qmc import Sobol, Halton, LatinHypercube

def pi_from(U):
    return 4.0 * ((U**2).sum(1) <= 1).mean()

Rs      = [2**k for k in range(5, 15)]   # 32 ... 16384
mc_err  = []
sob_err = []
hal_err = []

for R in Rs:
    rng = np.random.default_rng(0)
    mc_err.append( abs(pi_from(rng.random((R, 2))) - np.pi))
    sob_err.append(abs(pi_from(Sobol(2, scramble=True, seed=0).random(R)) - np.pi))
    hal_err.append(abs(pi_from(Halton(2, scramble=True, seed=0).random(R)) - np.pi))

plt.loglog(Rs, mc_err,  'b-o', label='MC (PCG64)')
plt.loglog(Rs, sob_err, 'r-s', label="Sobol'")
plt.loglog(Rs, hal_err, 'g-^', label='Halton')
ref = [Rs[0]**0.5 / r**0.5 * mc_err[0] for r in Rs]
plt.loglog(Rs, ref, 'k--', label=r'$O(R^{-1/2})$')
plt.xlabel('R'); plt.ylabel('|error|')
plt.grid(True); plt.legend()
plt.title('MC vs QMC convergence to $\pi$')
plt.tight_layout(); plt.show()
\end{lstlisting}

\end{frame}

%% ==================================================================
%%  SECTION 1.4
%% ==================================================================
\begin{frame}[allowframebreaks]{1.4 MC vs Quasi-MC: Comparison Examples}

  \begin{block}{Remark 1.4.2 -- Discrepancy in multidimensions}
    Grouping consecutive numbers from a 1-D low-discrepancy sequence
    into $d$-dimensional vectors \textbf{does not} preserve the
    low-discrepancy property.  Each dimension of the $d$-vector must be
    constructed simultaneously (as done by \texttt{scipy.stats.qmc}).
  \end{block}

  \textbf{Estimating $\pi$ (book p.~10):}
  With $R = 2^{10}$ Sobol' points the error is typically
  $10\text{--}50\times$ smaller than MC. The advantage is largest for
  smooth integrands and moderate $d$.

  \framebreak

  \begin{block}{Example 1.4.2 -- Winning Probability in a 2-D Game}
    \small
    Grid: $\{0,\ldots,N_1\}\times\{0,\ldots,N_2\}$. Starting from
    $(i_1,i_2)$:
    \begin{itemize}
      \item Any coordinate hits $0$ $\Rightarrow$ \textbf{LOSE}.
      \item Both coordinates reach $N_j$ $\Rightarrow$ \textbf{WIN}.
    \end{itemize}
    Transitions from $(i_1,i_2)$:
    \[
      (i_1,i_2)\to
      \begin{cases}
        (i_1+1,i_2) & \text{prob. } p_1(i_1),\\
        (i_1-1,i_2) & \text{prob. } q_1(i_1),\\
        (i_1,i_2+1) & \text{prob. } p_2(i_2),\\
        (i_1,i_2-1) & \text{prob. } q_2(i_1),\\
        (i_2,i_2)   & \text{prob. } 1-(p_1+q_1+p_2+q_2).
      \end{cases}
    \]
  \end{block}

  \framebreak

  \begin{block}{Exact winning probability (Formula 1.4.1)}
    \[
      \rho(i_1, i_2) \;=\;
      \frac{\displaystyle
        \prod_{j=1}^{2}\!\left(
          \sum_{n_j=1}^{i_j}\prod_{r=1}^{n_j-1}
          \tfrac{q_j(r)}{p_j(r)}
        \right)}
      {\displaystyle
        \prod_{j=1}^{2}\!\left(
          \sum_{n_j=1}^{N_j}\prod_{r=1}^{n_j-1}
          \tfrac{q_j(r)}{p_j(r)}
        \right)},
    \]
    where empty products equal 1.
  \end{block}

  \textbf{Symmetric example:}
  $N_1 = N_2 = 4$,
  $q_j(i_j) = p_j(i_j) = \tfrac{1}{4}$ for all $j, i_j$.
  Formula gives:
  \[
    \rho(2,3) = \frac{3}{8} = 0.375.
  \]

  \framebreak

  \textbf{Simulation algorithm -- one game from $(i_1,i_2)$:}

  Use two independent uniforms $U_1, U_2 \in [0,1)$ at each step:
  \begin{enumerate}
    \item Decide \emph{which coordinate} changes with $U_1$:
      \begin{itemize}
        \item Change coord.~1 if $U_1 \in [0,\ p_1(i_1)+q_1(i_1))$;
        \item Change coord.~2 if $U_1 \in [p_1+q_1,\ p_1+q_1+p_2+q_2)$;
        \item No change otherwise.
      \end{itemize}
    \item Decide \emph{direction} with $U_2$:
      \begin{itemize}
        \item Coord.~1 up if $U_2 < p_1(i_1)/(p_1(i_1)+q_1(i_1))$, else down.
        \item Coord.~2 up if $U_2 < p_2(i_2)/(p_2(i_2)+q_2(i_1))$, else down.
      \end{itemize}
    \item Stop if any coord hits 0 (LOSE) or all coords reach $N_j$ (WIN).
  \end{enumerate}

  \textbf{Estimator:}
  $\hat{Y}_R^{A_i} = \frac{1}{R}\sum_{j=1}^{R} Y_j^{A_i}$,
  where $Y_j^{A_i}\in\{0,1\}$ and
  $A_1=\text{PCG64}$, $A_2=\text{LatHyp}$, $A_3=\text{Sobol}$,
  $A_4=\text{Halton}$.

  \framebreak

  \begin{block}{Book Table~1.2 -- estimates of $\rho(2,3) = 0.375$}
    \centering
    \begin{tabular}{rcccc}
      \toprule
      Avg.\ steps & PCG64 & LatHyp & Sobol & Halton \\
      \midrule
      1\,000 & 0.3760 & 0.3686 & 0.3715 & 0.3750 \\
      5\,000 & 0.3748 & 0.3762 & 0.3748 & 0.3750 \\
      \bottomrule
    \end{tabular}
  \end{block}

  \textbf{Observations (Book Fig.~1.7):}
  \begin{itemize}
    \item All four generators converge to $0.375$.
    \item Sobol and Halton are more accurate for small $R$.
    \item The error curves for QMC are smoother (less variance).
  \end{itemize}

  \textbf{Structured points -- Book Fig.~1.9:}
  Scatter-plots of $(U_1^{(i)}, U_1^{(i+1)})$ for Sobol/Halton reveal
  clear \emph{linear structure} (``stratification lines''), unlike the
  uniform scatter of PCG64. This can cause artifacts for integrands
  aligned with these patterns.

  \begin{alertblock}{Practical rule of thumb}
    Use \textbf{QMC} for smooth, low-dimensional integrands.
    For $d \gtrsim 10$, non-smooth $f$, or Markov-chain-based algorithms,
    use plain \textbf{MC}.
  \end{alertblock}

\end{frame}

\begin{frame}[fragile]{1.4 Python: 2-D Game Simulation -- MC vs QMC}

\begin{lstlisting}[style=python, caption={Winning probability: MC vs Sobol estimation}]
import numpy as np
from scipy.stats.qmc import Sobol

N1, N2, p_val = 4, 4, 0.25

def play_one_game(stream):
    """stream: iterator yielding (U1, U2) pairs."""
    i1, i2 = 2, 3
    for U1, U2 in stream:
        thr1 = 2 * p_val                 # p1 + q1
        thr2 = thr1 + 2 * p_val          # + p2 + q2
        if U1 < thr1:
            i1 += 1 if U2 < 0.5 else -1
        elif U1 < thr2:
            i2 += 1 if U2 < 0.5 else -1
        if i1 == 0 or i2 == 0:  return 0
        if i1 == N1 and i2 == N2: return 1
    return 0

def mc_estimate(R, seed=0):
    rng  = np.random.default_rng(seed)
    wins = sum(play_one_game(iter(rng.random((500, 2)))) for _ in range(R))
    return wins / R

print(f"MC  (R=2000): {mc_estimate(2000):.4f}")
print(f"True value : 0.3750")
\end{lstlisting}

\end{frame}

%% ==================================================================
%%  SECTION 1.5
%% ==================================================================
\begin{frame}[allowframebreaks]{1.5 Bibliographical Notes}

  \textbf{Origins of Monte Carlo:}
  \begin{itemize}
    \item Developed at Los Alamos in the 1940s by von~Neumann, Ulam,
      and Metropolis for nuclear weapons simulation.
    \item Named after the Monte Carlo casino (Metropolis \& Ulam, 1949).
  \end{itemize}

  \textbf{Pseudorandom number generation:}
  \begin{itemize}
    \item \textbf{Knuth [75]}, Vol.~2, Ch.~3: comprehensive classical reference.
    \item L'Ecuyer [82]; Niederreiter [101].
  \end{itemize}

  \textbf{Random walks and arcsine law:}
  \begin{itemize}
    \item \textbf{Feller [52]}, Chapters~III, XIV: arcsine law, ballot problem.
    \item \textbf{Durrett [43]}: modern probability, random walk results.
    \item Law of the iterated logarithm: [100].
  \end{itemize}

  \textbf{Quasi-Monte Carlo:}
  \begin{itemize}
    \item \textbf{Niederreiter [101]}: \textit{Random Number Generation and
      Quasi-Monte Carlo Methods} --- foundational monograph.
    \item Halton [60]; Sobol' [119]; Faure [50].
    \item Tezuka [124]: scrambling and randomisation.
    \item \textbf{Glasserman [55]}: QMC in finance.
    \item Python: \texttt{scipy.stats.qmc} (SciPy $\ge$ 1.7).
  \end{itemize}

  \textbf{Travelling Salesman Problem:}
  \begin{itemize}
    \item \textbf{Applegate et al.\ [4]}: \textit{The Traveling Salesman Problem}.
    \item Simulated annealing: Kirkpatrick, Gelatt \& Vecchi [74] (1983).
  \end{itemize}

\end{frame}

%% ==================================================================
%%  SUMMARY
%% ==================================================================
\begin{frame}[allowframebreaks]{Chapter 1 Summary}

  \begin{block}{Core concepts introduced}
    \begin{enumerate}
      \item \textbf{Monte Carlo estimator}:
        $\hat{I}_R = \frac{1}{R}\sum_{i=1}^R f(U_i)$,
        RMSE $= \mathcal{O}(R^{-1/2})$, \emph{dimension-free}.

      \item \textbf{SSRW} $S_n = \sum_{j=1}^n X_j$:
        \begin{itemize}
          \item \emph{LIL}: fluctuations scale as $\pm\sqrt{2n\log\log n}$.
          \item \emph{Arcsine law}: $L^+_{2n}/(2n)$ has U-shaped arcsine
            distribution, not uniform on $[0,1]$.
        \end{itemize}

      \item \textbf{TSP}: NP-hard; MC heuristics (SA, CE) find
        near-optimal tours efficiently.

      \item \textbf{Quasi-MC}: low-discrepancy sequences (Sobol', Halton,
        golden-ratio) achieve $\mathcal{O}((\log R)^d/R)$ error for smooth $f$.

      \item \textbf{MC vs QMC}: QMC wins for smooth, low-$d$ problems;
        MC is more robust for high-$d$, discontinuous, or Markov-chain
        applications.
    \end{enumerate}
  \end{block}

  \framebreak

  \textbf{Key Python tools:}
  \begin{itemize}
    \item \texttt{numpy.random.default\_rng(seed)} --- reproducible PRNG
    \item \texttt{rng.random(n)}, \texttt{rng.integers(0, 2, n)} --- sampling
    \item \texttt{scipy.stats.qmc.Sobol}, \texttt{.Halton},
      \texttt{.LatinHypercube} --- QMC sequences
  \end{itemize}

  \bigskip
  \textbf{Looking ahead:}
  \begin{itemize}
    \item \textbf{Chapter~2:} How are pseudorandom numbers constructed,
      and how do we test their quality?
    \item \textbf{Chapter~3:} How to sample from arbitrary distributions
      (not just Uniform)?
    \item \textbf{Chapter~4:} Rigorous statistical analysis of simulation
      output -- confidence intervals and error estimation.
  \end{itemize}

\end{frame}

\end{document}
